\item El peso verdadero de un objeto se puede medir en un vacío, donde las fuerzas de flotación están ausentes. Un objeto de volumen $V$ se pesa en el aire sobre una balanza de brazos iguales, con el uso de contrapesos de densidad $\rho_c$. Al representar la densidad del aire como $\rho_{\text{aire}}$ y la lectura de la balanza como $F'_g$, demuestre que el verdadero peso $F_g$ es
$$ F_g = F'_g + \left( V - \frac{F'_g}{\rho_c g} \right) \rho_{\text{aire}} g $$
